\section{Background: From Script to Flow}
\label{sec:background}

\textbf{The grounding gap.} EDA scripting benchmarks (VerilogEval, RTLLM) evaluate isolated code generation: given a spec, produce Verilog. The agent makes 1--3 tool calls, never inspects a live tool's output, and never recovers from an error. \bench{} introduces \emph{design-flow orchestration}: the agent must produce a deliverable (\texttt{submit.zip}) by orchestrating remote commercial EDA tools through an MCP bridge, scored against $\geq 8$ executable assertions. The median run issues 26.5 tool calls (IQR 13.8--38.0) across 20.5 iterations, with 41.2\% of calls going through the MCP bridge to live EDA tools (47 instrumented pilot traces; \S\ref{sec:analysis}). This is qualitatively harder: the agent must ground its actions in tool outputs it has never seen during training, chain dependent steps, and recover from failures mid-flow.

\textbf{Why memory and RSI may matter more in flow.} On single-script tasks, in-context learning within one trace often suffices---the entire problem fits in the window. On long-chain flow, the relevant knowledge (tool quirks, project conventions, failure patterns) accumulates \emph{across} sessions and exceeds any window. This makes persistent memory a potential binding constraint: an agent that can consolidate ``which tool sequences work'' and ``which errors recur'' into a cross-session memory may close the grounding gap, while one that re-derives everything from scratch each run is bounded. Recursive self-improvement (RSI)---expanding a skill library and refining the harness from experience---is only meaningful if memory persists to accumulate that experience. Our prior work on script-layer EDA (DATE27, \cite{PLACEHOLDER_DATE27_ZHULONG}; ZhuLong-DAC, \cite{liu2026runreadbackrefine}) showed memory + RSI gains on short-chain tasks; whether these \emph{transfer} to long-chain flow is the central question of this paper.

\textbf{From biological affect to computational mapping.} Cognitive neuroscience establishes that emotional arousal modulates memory consolidation (amygdala--hippocampus interaction; \cite{mcgaugh2004}), that high arousal triggers a central/peripheral encoding trade-off (weapon-focus; \cite{loftus1987}), and that retrieval is mood-congruent and state-dependent \cite{bower1981}. We map these three principles to computable mechanisms for an EDA agent's memory architecture (\S\ref{sec:harness}), turning ``does memory help?'' into ``\emph{which} mechanism brings the gain, and under what conditions does it vanish?''
